4) let $G$ be a group. Take $A = G$
Define $*: G \times A \to A$ as
$g * a = gag^{-1} \quad \forall g \in G, \forall a \in A$
Then $*$ is a group action on $G$ and is known as action by conjugation.
Take $g_1, g_2 \in G$ and $a \in A(=G)$
$g_1 * (g_2 * a) = g_1 * (g_2 a g_2^{-1})$
$= g_1 g_2 a g_2^{-1} g_1^{-1}$
$= g_1 g_2 a (g_1 g_2)^{-1}$
$g_1 * (g_2 * a) = (g_1 g_2) * a \longrightarrow (1)$
$e * a = eae^{-1}$
$e * a = a \longrightarrow (2)$

5) let $G$ be a group and $H \trianglelefteq G$. Then take
$A = G/H = \text{set of all left coset of } H \text{ in } G$.
i.e. $A = \{aH \mid a \in G\}$
Define, $*: G \times A \to A$ by
$g * (xH) = gxg^{-1} H \quad \forall g \in H$
Take $g_1, g_2 \in G \ \& \ xH \in A$
Then $g_1 * (g_2 * xH) = g_1 * (g_2 x g_2^{-1} H)$
$= g_1 g_2 x g_2^{-1} g_1^{-1} H$
$= g_1 g_2 x (g_1 g_2)^{-1} H$
$g_1 * (g_2 * xH) = (g_1 g_2) * xH$
$e * xH = exe^{-1} H$
$= xH$
$\therefore *$ is a group action.
